\documentclass[11pt,a4paper]{article}

\usepackage[utf8]{inputenc}

\usepackage[T1]{fontenc}

\usepackage[polish]{babel}

\usepackage{lmodern}

\usepackage{geometry}

\geometry{margin=27mm}

\usepackage[hidelinks]{hyperref}

\title{Indywidualny dobór strategii regulacji i skuteczność w kontekście}

\author{Mateusz Pomianek}

\date{Szkic naukowy, wersja 1.0.0 | 28 listopada 2026}

\begin{document}

\maketitle

\begin{abstract}

Proponowana miara główna to różnica efektów przydziału strategii między kontekstami na skali prawdopodobieństwa poprawnego powrotu.

\end{abstract}

\section{Wprowadzenie}

Elastyczność regulacyjna zakłada dopasowanie strategii do dostępnych możliwości działania (Bonanno \& Burton, 2013). Mikrorandomizacja pozwala szacować efekt bliski momentowi przydziału strategii (Klasnja et al., 2015).

\section{Cel i pytania badawcze}

Hipotezy robocze wymagają zatwierdzenia przed rozpoczęciem właściwej weryfikacji.

\paragraph{H1.} Podpowiedź ukierunkowana na zmianę działania zwiększy prawdopodobieństwo poprawnego powrotu w ciągu 30 s względem krótkiej pauzy, gdy zadanie oferuje dostępny sposób naprawienia błędu. Przewaga będzie mniejsza w sytuacji, w której naprawienie skutku błędu jest niedostępne.

\section{Materiał i metody}

Pierwszorzędowy wskaźnik: różnica efektów przydziału strategii między kontekstami na skali prawdopodobieństwa poprawnego powrotu.

W dostępnych momentach randomizacji należy przydzielić jedną z dwóch podpowiedzi i modelować interakcję przydziału z ustalonym przed randomizacją kontekstem. Zmienność efektów między osobami wymaga osobnej analizy ze skurczem oszacowań. Należy ustalić reguły kolejnej randomizacji i odstęp między podpowiedziami.

Przed badaniem z udziałem ludzi należy zatwierdzić wymagania etyczne, zasady zgody i ochronę danych.

\section{Plan analizy}

Zamrożenie estymandu, reguł oceny jakości i obsługi braków nastąpi przed dostępem do danych wynikowych. Raport obejmie oszacowanie kontrastu głównego z przedziałem niepewności i pełne mianowniki. Jeśli przewidziano hipotezy drugorzędowe, przed analizą należy ustalić sposób kontroli wielokrotności.

\section{Status dowodów i granice własności}

Koncepcja figuruje w katalogu programu ARC Neuro jako MSC-ADHD-P07 (IDEA/DEFER). Publiczne rejestry \texttt{config/public-projects.json} nie wskazują jeszcze kanonicznego projektu ani repozytorium właściciela tego manuskryptu. Autorstwo szkicu nie jest potwierdzeniem ustalonej własności repozytoryjnej. Zgodnie z architekturą programu właściwy projekt badawczy powinien przechowywać zatwierdzony manuskrypt, protokół, dane i rejestr dowodów. Ten plik nie stanowi ich kopii.

Zamrożenie estymandu i reguł analizy musi nastąpić przed dostępem do danych.

\section{Wyniki}

Nie zebrano danych empirycznych dla tej koncepcji.

\section{Dyskusja i ograniczenia}

Ograniczenia obejmują możliwe błędy pomiaru, dobór próby, alternatywne mechanizmy i ograniczoną możliwość uogólniania. Korelacja nie dowodzi przyczynowości; poprawność techniczna nie jest dowodem znaczenia klinicznego.

\section{Wnioski}

Hipotezy i sposób pomiaru stanowią podstawę oceny wykonalności oraz przygotowania właściwego protokołu.

\section*{Bibliografia}

{\setlength{\parindent}{0pt}
\setlength{\parskip}{5pt}
\hangindent=1.27cm
Bonanno, G. A., \& Burton, C. L. (2013). Regulatory flexibility: An individual differences perspective on coping and emotion regulation. \textit{Perspectives on Psychological Science, 8}(6), 591--612. \url{https://doi.org/10.1177/1745691613504116}\par
\hangindent=1.27cm
Klasnja, P., Hekler, E. B., Shiffman, S., Boruvka, A., Almirall, D., Tewari, A., \& Murphy, S. A. (2015). Microrandomized trials: An experimental design for developing just-in-time adaptive interventions. \textit{Health Psychology, 34}(Suppl.), 1220--1228. \url{https://doi.org/10.1037/hea0000305}\par
}

\end{document}
